\chapter{Doce casos de teoría avanzada para examen}
\label{ch:matmo-teoria-avanzada}
\noindent Cuestionario original, separado de las preguntas de los simulacros S01--S10. Cada caso vale 5 puntos académicos: 2 por definición/mecanismo, 2 por argumentación y 1 por límites o datos que faltan. No representa el formato oficial de la ULE.

\section*{AT01. Fractura y geometría}
Una placa presenta una grieta central total de 8 mm. Explique qué representa $a$ en $K_I=Ysigmasqrt{pi a}$, por qué no puede usar el mismo $Y$ para una grieta de borde sin justificarlo y qué validez de tenacidad debe comprobar.

\section*{AT02. Fatiga y tensión media}
Dos componentes presentan igual $sigma_{max}$, pero uno tiene $R=-1$ y el otro $R=0$. ¿Por qué no tienen necesariamente la misma vida? Razone qué se debe calcular antes de consultar una curva S--N.

\section*{AT03. Paris y predicción de inspecciones}
Un estudiante integra Paris entre dos longitudes y llama al resultado intervalo de inspección permitido. Identifique al menos tres saltos lógicos.

\section*{AT04. Fluencia y álabes monocristalinos}
Compare fluencia por movimiento de dislocaciones y mecanismos asociados a límites de grano, y explique el objetivo de una estructura monocristalina.

\section*{AT05. Material compuesto orientado}
Se proponen dos módulos calculados, $E_L$ y $E_T$, diferentes para una lámina de fibra. ¿Qué hipótesis y qué mecanismos explican la diferencia? ¿Por qué no se debe tratar el laminado como isotrópico?

\section*{AT06. Selección END}
Una grieta cerrada subsuperficial en material no ferromagnético debe inspeccionarse. ¿Por qué penetrantes y partículas magnéticas no se eligen directamente? ¿Qué necesita saber para comparar UT y radiografía?

\section*{AT07. Diseño, aeronave y mantenimiento}
Distinga certificado de tipo, certificado individual de aeronavegabilidad y CRS tras mantenimiento. ¿Puede cualquiera de los tres sustituir los otros dos?

\section*{AT08. Parte 145, 66 y 147}
Un técnico supera formación Parte 147 y posee licencia Parte 66, pero trabaja fuera del alcance autorizado de su organización. ¿Puede emitir un CRS por ello?

\section*{AT09. AD y boletín de servicio}
Una AD cita la revisión 1 de un SB y el fabricante emite la revisión 3. ¿Qué debe verificarse antes de usar la revisión 3 para satisfacer la AD?

\section*{AT10. ARC con hallazgos y plazos}
La revisión detecta un hallazgo abierto y se pretende expedir ARC porque restan 20 días hasta caducidad. Explique el error. Diferencie adelanto de revisión física y extensión del ARC.

\section*{AT11. Tarea con dos límites y LLP}
Una tarea vence por FH o fecha, lo que ocurra primero, mientras una LLP presenta ciclos totales desconocidos. ¿Qué puede decidirse con datos fiables y qué queda indeterminado?

\section*{AT12. EGT, vibración y diagnóstico}
Después de una limpieza mejora el margen EGT, pero persiste vibración de primer orden y hay partículas en aceite. ¿Puede darse por resuelto el problema?
